\noindent\textbf{Source:} M. Giaquinta, \emph{Introduction to Regularity Theory for Nonlinear Elliptic Systems}.\\
\textbf{Topic:} Elliptic regularity / Energy estimates.\\
\textbf{Status:} Personal mathematical exploration.
\section{Caccioppoli Inequality and the Role of Cutoff Functions}

\subsection*{Motivation}
In the standard proof of the Caccioppoli inequality (see, for instance, M. Giaquinta), a cutoff function $\eta$ is introduced as a technical device for localising the energy estimate. The cutoff is typically chosen to satisfy $\eta \equiv 1$ on the inner ball $B_\rho$, $\eta \equiv 0$ outside the outer ball $B_R$, and $|\nabla \eta| \le C/(R-\rho)$.

This raised a question in my mind: is the assumption $\eta \equiv 1$ on the inner ball merely a convenient choice, or is it fundamentally necessary for the classical unweighted estimate? To explore this, I decided to modify the cutoff function so that it varies inside the inner ball, and examined what happens to the inequality.

\begin{quote}
    How does the classical Caccioppoli inequality change if the cutoff function is not identically equal to one and varies inside the inner ball?
\end{quote}

\begin{question}
    Let $u \in W^{1,2}(\Omega)$ be a weak solution of $\Delta u = 0$. What happens to the classical Caccioppoli inequality
    \[
    \int_{B_\rho} |\nabla u|^2 \le \frac{C}{(R-\rho)^2} \int_{B_R} |u|^2
    \]
    if the cutoff function varies inside the inner ball and is never identically equal to 1 in this region?
\end{question}

\begin{answer}
    The classical Caccioppoli inequality requires the cutoff function $\eta \equiv 1$ on the inner ball $B_\rho$. The gradient of the cutoff function satisfies
    \begin{equation}
        |\nabla \eta| \le \frac{C}{R-\rho},
    \end{equation}
    where $C > 0$ is a constant.

    Consider the test function $\phi = \eta^2 u \in W^{1,2}_0(B_R)$, as in the standard proof. Since $\Delta u = 0$, multiplying by $\phi$ and integrating over the domain gives
    \[
    \int_{\Omega} \Delta u \, \phi = 0.
    \]
    Using integration by parts,
    \[
    -\int_\Omega \nabla u \cdot \nabla \phi = 0,
    \]
    or equivalently,
    \begin{equation}
        \int_\Omega \nabla u \cdot \nabla \phi = 0. \tag{1}
    \end{equation}

    Differentiating $\phi = \eta^2 u$ gives
    \[
    \nabla \phi = \eta^2 \nabla u + 2\eta u \nabla \eta.
    \]
    Substituting this into (1) yields
    \[
    \int_\Omega \eta^2 |\nabla u|^2 + 2 \int_\Omega \eta u \nabla u \cdot \nabla \eta = 0.
    \]
    Therefore,
    \begin{align*}
        \int_\Omega \eta^2 |\nabla u|^2 &= -2 \int_\Omega \eta u \nabla u \cdot \nabla \eta \\
        &\le 2 \int_\Omega \eta |u| |\nabla u| |\nabla \eta|.
    \end{align*}
    Restricting the domain of integration to the outer ball $B_R$, we obtain
    \begin{equation}
        \int_{B_R} \eta^2 |\nabla u|^2 \le 2 \int_{B_R} \eta |u| |\nabla u| |\nabla \eta|. \tag{2}
    \end{equation}

    Applying Hölder's inequality to the right-hand side of (2),
    \begin{equation}
        \int_{B_R} \eta^2 |\nabla u|^2
        \le
        \left( \int_{B_R} \eta^2 |\nabla u|^2 \right)^{1/2}
        \left( \int_{B_R} 4 |u|^2 |\nabla \eta|^2 \right)^{1/2}. \tag{3}
    \end{equation}
    Dividing both sides of (3) by $\left( \int_{B_R} \eta^2 |\nabla u|^2 \right)^{1/2}$ and then squaring gives
    \begin{equation}
        \int_{B_R} \eta^2 |\nabla u|^2 \le \int_{B_R} 4 |u|^2 |\nabla \eta|^2. \tag{4}
    \end{equation}
    Using the bound (1) for $|\nabla \eta|$, we obtain
    \begin{equation}
        \int_{B_R} \eta^2 |\nabla u|^2 \le \int_{B_R} \frac{C}{(R-\rho)^2} |u|^2, \tag{5}
    \end{equation}
    where $C = 4C^2$ is a new constant.

    Since the inner ball is contained in the outer ball, we have
    \[
    \int_{B_\rho} \eta^2 |\nabla u|^2 \le \int_{B_R} \eta^2 |\nabla u|^2,
    \]
    and hence
    \begin{equation}
        \int_{B_\rho} \eta^2 |\nabla u|^2 \le \int_{B_R} \frac{C}{(R-\rho)^2} |u|^2. \tag{6}
    \end{equation}
    That is,
    \begin{equation}
        \boxed{
            \int_{B_\rho} |\eta \nabla u|^2
            \le
            \frac{C}{(R-\rho)^2}
            \int_{B_R} |u|^2
        }. \tag{7}
    \end{equation}

    Thus, when the cutoff function varies inside the inner ball and is not identically equal to 1, the resulting inequality differs from the classical Caccioppoli inequality. The difference lies in the left-hand side: the energy of the gradient is replaced by the energy of the weighted gradient $\eta \nabla u$.

    \subsection*{Further Thoughts}
    This observation clarifies the role played by the cutoff function in the classical proof. The gradient energy on the inner ball is controlled by the weighted gradient energy, where the weight is the cutoff function $\eta$. This modified inequality controls only the weighted gradient on the inner ball, not the full gradient. If $\eta \equiv 1$ on $B_\rho$, the classical inequality is recovered. This shows that the condition $\eta \equiv 1$ on the inner ball is not merely a convenience — it is essential for obtaining the unweighted gradient estimate.
\end{answer}